**Title**  
Exploring Multimodal Adaptation Frameworks for Emotion Recognition: A Case Study on Indonesian Cultural Contexts  

---

**Abstract**  
Emotion recognition research in low-resource and culturally specific languages remains underexplored despite its critical importance for social media and interactive systems. This study introduces OmniMER, a multimodal adaptation framework, and benchmarks it on IndoMER, the first multimodal Indonesian emotion recognition dataset. IndoMER consists of 1,944 annotated video segments from 203 speakers across text, audio, and visual modalities, reflecting authentic cultural challenges such as cross-modal inconsistency and class imbalance. OmniMER enhances emotion recognition by incorporating auxiliary tasks for text, audio, and video perception before fusing modalities. Results demonstrate OmniMER’s significant improvements in Macro-F1 scores (7.6 and 22.1 points for sentiment and emotion classification, respectively) compared to baseline models. Cross-lingual testing on Chinese datasets further validates OmniMER’s generalizability. This paper contributes practical insights into multimodal emotion recognition frameworks for culturally nuanced data and presents IndoMER as a benchmark for future research.  

---

**Introduction**  
Emotion recognition plays a pivotal role in interactive systems, social media analysis, and mental health monitoring. While large datasets have enabled progress in languages like English and Chinese, underserved languages such as Indonesian remain under-researched, despite their prominence in Southeast Asia. Indonesian cultural norms, characterized by indirect communication and subtle emotional expression, pose unique challenges to multimodal emotion recognition. To address these challenges, we introduce OmniMER, an auxiliary-enhanced multimodal adaptation framework. OmniMER builds upon large language models (LLMs) and leverages auxiliary tasks specific to text, audio, and visual modalities to improve emotion recognition in culturally complex datasets.  

This study benchmarks OmniMER on IndoMER, a new multimodal Indonesian emotion recognition dataset, and explores its generalizability across other languages. IndoMER captures seven emotion categories in video segments, reflecting real-world challenges such as cross-modal inconsistency and long-tailed class distributions. By integrating auxiliary tasks, OmniMER reduces reliance on spurious correlations and enhances interpretability in low-resource settings.  

---

**Related Work**  
Emotion recognition research has primarily focused on resource-rich languages and datasets. Methods like Partial Information Decomposition (Ma et al., 2025) have improved interpretability in multimodal fusion, while frameworks such as VideoScaffold (Zheng et al., 2025) address temporal coherence in video-based models. However, these studies lack consideration of cultural nuances and language-specific characteristics. Recent advances in multimodal LLM adaptation, such as OmniMER, aim to bridge this gap by leveraging auxiliary tasks for modality-specific perception.  

In the cultural context, few studies explore datasets reflecting authentic communication norms. IndoMER fills this void, providing temporally aligned multimodal annotations for Indonesian speakers. The integration of auxiliary tasks for emotion keyword extraction, prosody analysis, and facial expression detection builds upon frameworks like Brain-Grounded Axes (Andric, 2025) and Explainable Multimodal Regression (Ma et al., 2025). This study extends these contributions by focusing on culturally nuanced emotion recognition and evaluating cross-lingual generalization.  

---

**Methods/Experiment**  
### Dataset  
IndoMER consists of 1,944 video segments annotated across text, audio, and visual modalities. Seven emotion categories are included: happiness, sadness, anger, surprise, disgust, fear, and neutral. The dataset captures cross-modal inconsistencies and long-tailed distributions, reflecting Indonesian communication styles.  

### Framework: OmniMER  
OmniMER incorporates three auxiliary tasks to enhance multimodal fusion:  
1. **Emotion Keyword Extraction**: Identifies emotion-relevant text features using attention mechanisms.  
2. **Facial Expression Analysis**: Detects visual cues using convolutional modules fine-tuned on facial landmarks.  
3. **Prosody Analysis**: Extracts audio signals linked to pitch, tone, and rhythm for emotional inference.  

These tasks are implemented through adaptive prefix-tuning and fused via a hierarchical transformer model.  

### Experimental Setup  
Experiments were conducted on IndoMER and CH-SIMS (Chinese dataset) to evaluate cross-lingual performance. Models were assessed using Macro-F1 scores for sentiment classification and emotion recognition. Baselines included standard LLMs and multimodal fusion strategies without auxiliary tasks.  

---

**Results**  
### Performance on IndoMER  
OmniMER significantly outperformed baseline models, achieving Macro-F1 scores of 0.582 for sentiment classification and 0.454 for emotion recognition. Improvements were observed across all modalities, with text contributing the most to classification accuracy.  

| **Metric**         | **Baseline** | **OmniMER** | **Improvement** |  
|---------------------|--------------|--------------|------------------|  
| Sentiment Macro-F1  | 0.506        | 0.582        | +7.6 points      |  
| Emotion Macro-F1    | 0.233        | 0.454        | +22.1 points     |  

### Cross-Lingual Evaluation  
Testing on CH-SIMS demonstrated OmniMER’s generalizability, with improvements of 15.4 points in Macro-F1 for emotion recognition compared to baseline models.  

| **Dataset** | **Metric** | **Baseline** | **OmniMER** | **Improvement** |  
|-------------|------------|--------------|--------------|------------------|  
| IndoMER     | Sentiment  | 0.506        | 0.582        | +7.6 points      |  
| CH-SIMS     | Emotion    | 0.300        | 0.454        | +15.4 points     |  

---

**Discussion**  
The results highlight OmniMER’s capacity to address challenges in multimodal emotion recognition for culturally specific datasets. The integration of auxiliary tasks improved consistency and interpretability by leveraging modality-specific cues. Notably, text-based emotion keyword extraction contributed most to sentiment classification, while visual and audio tasks enhanced emotion recognition.  

Cross-lingual evaluation revealed OmniMER’s adaptability, suggesting its potential application to other low-resource languages. However, challenges remain in optimizing auxiliary tasks for datasets with severe class imbalance and cross-modal inconsistencies. Future work will explore dynamic weighting mechanisms to balance modality contributions and refine auxiliary task architectures.  

---

**Conclusion**  
This study introduced OmniMER, a multimodal adaptation framework, and benchmarked it on IndoMER, the first Indonesian multimodal emotion recognition dataset. OmniMER’s auxiliary-enhanced architecture significantly improved Macro-F1 scores for sentiment and emotion classification, outperforming baseline models. Cross-lingual evaluation demonstrated its generalizability to other languages, providing a foundation for future research on culturally nuanced emotion recognition.  

---

**Contributions**  
1. Developed OmniMER, an auxiliary-enhanced multimodal framework for emotion recognition.  
2. Introduced IndoMER, a benchmark dataset reflecting Indonesian communication norms.  
3. Demonstrated OmniMER’s effectiveness in improving sentiment and emotion recognition accuracy.  
4. Validated OmniMER’s generalizability through cross-lingual testing on CH-SIMS.  
5. Provided insights into modality-specific contributions and future optimization strategies.  

---  
This paper contributes to the advancement of multimodal emotion recognition frameworks, emphasizing the importance of cultural context and auxiliary tasks for improving model performance in low-resource settings.